\section{Amplifying an average degree saving}\label{sec:mean-cost}

A finite observation can have some inexpensive outputs and some expensive
ones. We now convert an average degree saving into a bound valid on every
input. Apply the observation independently many times, retain output
words whose total cost is controlled, and merge all other words into one
label. The indicator of the merged cell is the complement of the retained
indicators. Its degree therefore satisfies the same bound, even if
individual discarded words do not.

This operation must also retain enough variance and restore the normal
approximation needed to repeat it. We use the score $T_F$, retained
variance $v(F)$, and pointwise cell-degree calculus of
Section~\ref{sec:observations}. The new ingredient is a separate cost for
each output label. These costs measure degrees of indicator expansions,
not the number of queries required to determine an output.

\subsection{From a variance saving to an average cost saving}
\label{sec:refinement-first}


Let $q\ge2$ be an integer, and let
$\mathcal P:\R^q\to\mathcal L$ be a measurable map to a finite label
set.  Give label $l$ an integer cost $\ell(l)\in\{1,\ldots,q\}$.
The relevant algebraic requirement is that, on every product of finite
input alphabets, the cell indicator for $l$ is a finite linear
combination of functions each depending on at most $\ell(l)$ input
coordinates.  The indicator itself need not depend on only that many
coordinates.  When the inputs are scores of observations with cell
degree bound $D$, this requirement bounds the degree of the output
cell by $D\ell(l)$.


Suppose now that a Gaussian construction supplies an event $A$ on which
one coordinate of degree cost can be saved, while the sum has conditional
variance less than one. The next lemma turns this event into an
observation whose retained variance exceeds its average cost. We keep
$A$ as one cell and refine its complement finely enough that the extra
prediction error uses less than the available variance saving.

\begin{lemma}[Refining the complement of one event]
\label{lem:refinement}
Let $q\ge2$, let $G=(G_1,\ldots,G_q)$ have independent standard normal
coordinates, and put $S_G=\sum_iG_i$.  Suppose a measurable set
$A\subset\R^q$ has Gaussian-null boundary,
\[
 p_A:=\Pp(G\in A)>0,\qquad
 r_A:=\Var(S_G\mid G\in A)<1,
\]
and its indicator, on every product of finite alphabets, is a finite
linear combination of functions each depending on at most $q-1$
coordinates.  Then there is a finite pointwise partition $\mathcal P$
with one cell $A$ of cost $q-1$ and all other cells of cost $q$ that
has Gaussian-null cell boundaries and the prescribed pointwise
coordinate bounds. Writing
\[
 V=\Var(\E[S_G\mid\mathcal P(G)]),\qquad
 c=\E\ell(\mathcal P(G)),
\]
we have $V>c=q-p_A$.
\end{lemma}

\begin{proof}
Partition $[-M,M)$ into finitely many half-open intervals of length
at most $h$, and add the two tail intervals.  Approximate a real
coordinate by the midpoint of its bounded interval and by zero on
the tails.  For a standard normal coordinate, the expected squared
error is at most
\[
 h^2+\E[G_1^2\one_{\{|G_1|\ge M\}}].
\]
Let $\widetilde S_G$ be the sum of these coordinate approximations.
It is constant on each box of the resulting finite rectangular grid.
Since $(\sum_i a_i)^2\le q\sum_i a_i^2$, choosing finite $M$ large
enough and then $h>0$ small enough ensures
\[
 \E(S_G-\widetilde S_G)^2<p_A(1-r_A).
\]

Give $A$ one label and give every intersection of $A^c$ with a grid
box its own label.  Retain this partition pointwise, even on cells
of zero Gaussian probability.  Every new cell has boundary contained
in the boundary of $A$ together with the finitely many grid
hyperplanes, and hence has Gaussian-null boundary.
Give $A$ cost $q-1$ and every other cell cost $q$.  The special-cell
algebraic requirement is assumed; the others are automatic for
functions of $q$ inputs.  Conditional means minimize squared error
on every cell~\cite[Theorem~4.1.15]{Durrett}, so
\[
 \E\Var(S_G\mid\mathcal P(G))
 \le p_A r_A+
       \E[\one_{A^c}(S_G-\widetilde S_G)^2]<p_A.
\]
Since $\Var(S_G)=q$, it follows that
$V>q-p_A=\E\ell(\mathcal P(G))=c$.
\end{proof}

The event is fixed before the complement grid is chosen. The resulting
partition is then fixed before any normal approximation or iteration.
This gives the concrete inequality $V>c$ that the amplification theorem
below needs.

\subsection{Amplifying the average cost saving}

\begin{theorem}[Amplification from a mean cost advantage]
\label{thm:mean-cost}
Let $q\ge2$, let $\mathcal P:\R^q\to\mathcal L$ be a measurable map
to a finite label set, and let $\ell:\mathcal L\to\{1,\ldots,q\}$.
Suppose that, on every product of finite input alphabets, each cell
indicator for $l$ is a finite linear combination of functions of at
most $\ell(l)$ input coordinates.  Suppose also that every cell has
boundary of standard Gaussian measure zero in $\R^q$.  For independent standard
normals $G_1,\ldots,G_q$, put
\[
 S_G=\sum_{i=1}^qG_i,\qquad
 V=\Var(\E[S_G\mid\mathcal P(G)]),\qquad
 c=\E\ell(\mathcal P(G)).
\]
If $V>c$, then finite observations on sign cubes have arbitrarily
large ratios $v(F)/D$, where $D>0$ is a cell degree bound.
The Boolean conclusion of Theorem~\ref{thm:main} follows by taking the
sign of the final score, without a further averaging step.
All partitions and observations are defined pointwise, including
cells of zero Gaussian probability.
\end{theorem}

Only restricted first moments need to pass to the Gaussian limit in
this criterion.  The next lemma isolates that fact; it also explains
why zero limiting cell probabilities cause no difficulty.

\begin{lemma}[Transfer for a fixed finite partition]
\label{lem:mc-transfer}
Under the boundary hypothesis of Theorem~\ref{thm:mean-cost}, let
$Z_n$ be centered random variables of variance one converging in
distribution to a standard normal.  Use $q$ independent copies to set
\[
 S_n=\sum_{i=1}^qZ_{n,i},\qquad L_n=\mathcal P(\mathbf Z_n),\qquad
 V_n=\Var(\E[S_n\mid L_n]),\qquad c_n=\E\ell(L_n).
\]
Then
\[
 \liminf_n V_n\ge V,\qquad c_n\longrightarrow c,
 \qquad \E|Z_n|\longrightarrow\sqrt{2/\pi}.
\]
No fourth-moment hypothesis is required.
\end{lemma}

\begin{proof}
Independence gives joint weak convergence of the input tuples.  The
Gaussian-null boundaries imply convergence of all cell probabilities
\cite[Theorem~3.10.1(v)]{Durrett}.
For a fixed cell $E$, clip the sum continuously to $[-t,t]$.
The clipped sum times $\one_E$ is bounded and has Gaussian-null
discontinuities, so its expectations converge
\cite[Theorem~3.10.1(vi)]{Durrett}.  Since $S_n$ is centered
and has variance $q$,
\[
 \E\bigl[|S_n|\one_{\{|S_n|>t\}}\bigr]\le q/t.
\]
Removing the clipping therefore proves
$\E[S_n\one_{\{\mathbf Z_n\in E\}}]\to
\E[S_G\one_{\{G\in E\}}]$.

For each $n$, the centered conditional mean satisfies
\[
 V_n=\sum_{l:\Pp(L_n=l)>0}
 \frac{\E[S_n\one_{\{L_n=l\}}]^2}{\Pp(L_n=l)}.
\]
Every term for a cell of positive Gaussian probability converges to
the corresponding Gaussian term.  All remaining terms are
nonnegative.  This proves the lower limit inequality; it does not
require deleting the cells with zero limiting probability.
The cost convergence follows from the finitely many convergent cell
probabilities.  Finally, weak convergence and the bound
$\E[|Z_n|\one_{\{|Z_n|>t\}}]\le1/t$ give convergence of the
absolute first moments.
\end{proof}

\begin{proof}[Proof of Theorem~\ref{thm:mean-cost}]
The partition and costs are fixed throughout the proof.  Choose
\[
 0<\delta<(V-c)/2,
 \qquad 1<\lambda<\frac{V-\delta}{c+\delta}.
\]
For any centered, variance-one law $Z$, use independent copies to
define
\[
 S_Z=\sum_{i=1}^q Z_i,\qquad L=\mathcal P(\mathbf Z),\qquad
 V_Z=\Var(\E[S_Z\mid L]),\qquad c_Z=\E\ell(L).
\]
For a finite observation $F$ with $v=v(F)>0$, write $Z=T_F/\sqrt v$.
We will increase $v/D$ while maintaining
\begin{equation}\label{eq:invariant}
 V_Z>V-\delta,\qquad c_Z<c+\delta/2,\qquad \E|Z|>1/2.
\end{equation}
Lemma~\ref{lem:mc-transfer} shows that all three inequalities hold
eventually along every sequence of centered, variance-one laws
converging to a standard normal.

To initialize, observe the bit sum itself: $F=H_N$, with score
$T_F=H_N$ and $v=D=N$.  The central limit theorem for independent
uniform signs~\cite[Theorem~3.4.1]{Durrett} supplies a finite $N$ for which
\eqref{eq:invariant} holds.  The initial ratio is $v/D=1$.

Now hold one such $F$ fixed.  Apply $\mathcal P$ to the standardized
scores of $q$ independent copies of $F$, obtaining the group label $L$.
The algebraic assumption and the degree calculus give
\[
 \deg(\one_{\{L=l\}})\le D\ell(l).
\]
The tower property identifies the group score and its variance as
\begin{equation}\label{eq:group}
 T_L=\sqrt v\,\E[S_Z\mid L],\qquad
 \tau^2:=v(L)=vV_Z>0.
\end{equation}
Thus this group has a variance advantage relative to its mean cost.
We next obtain a common degree bound without losing that advantage.

Take $K$ independent copies of $L$.  Retain a label word
$(l_1,\ldots,l_K)$ when $\sum_{b=1}^K\ell(l_b)\le K(c+\delta)$.
Let $F'_K$ report each retained word exactly and report one failure
label for every other word.  Every retained-word indicator is a product
with degree at most
\begin{equation}\label{eq:degree}
 D'_K=DK(c+\delta).
\end{equation}
The failure indicator is one minus the finite sum of all retained-word
indicators, so it satisfies the same degree bound on every input.
The real upper bound $D'_K$ need not be rounded to an integer.
Since the independent costs have mean $c_Z<c+\delta/2$, the law of
large numbers~\cite[Theorem~2.2.3]{Durrett} gives a failure probability
$\varepsilon_K\to0$.

Write $U_K=\sum_{b=1}^KT_{L_b}$ and $T'_K=T_{F'_K}$.
The tower property again gives $T'_K=\E[U_K\mid F'_K]$.
On a retained word, $U_K$ is known exactly.  On failure, its conditional
mean is at least as good a predictor as zero
\cite[Theorem~4.1.15]{Durrett}.  Consequently
\begin{align}
 \E(U_K-T'_K)^2
 &\le\E[U_K^2\one_{\{\mathrm{failure}\}}]\notag\\
 &\le(\E U_K^4)^{1/2}\varepsilon_K^{1/2}
   =o(K\tau^2). \label{eq:loss}
\end{align}
For the last equality, $T_L$ has finite range, mean zero, and positive
variance, and independence gives
\[
 \E U_K^4=K\E T_L^4+3K(K-1)\tau^4=O(K^2\tau^4).
\]
The implicit constant can depend on the current observation, which
has been fixed before $K$ tends to infinity.

The conditional-expectation projection identity
\cite[Theorem~4.1.15 and its proof]{Durrett} now yields
\[
 v(F'_K)=K\tau^2-\E(U_K-T'_K)^2\sim K\tau^2.
\]
In particular, this variance is positive for all sufficiently large
$K$.  The central limit theorem~\cite[Theorem~3.4.1]{Durrett} gives
$U_K/(\sqrt K\tau)\Rightarrow N(0,1)$.  By \eqref{eq:loss}, replacing
$U_K$ by $T'_K$ changes the normalized variable by a term tending to
zero in $L^2$.  Moreover,
$\sqrt{v(F'_K)}/(\sqrt K\tau)\to1$.  The addition and rescaling forms of
Slutsky's theorem~\cite[Exercises~3.2.13--3.2.14]{Durrett} therefore give
\[
 \frac{T'_K}{\sqrt{v(F'_K)}}\Rightarrow N(0,1).
\]
These new standardized scores are centered and have variance one.
Lemma~\ref{lem:mc-transfer} restores all three inequalities in
\eqref{eq:invariant} for sufficiently large $K$.  At the same time,
\eqref{eq:group} and \eqref{eq:degree} give
\[
 \frac{v(F'_K)}{D'_K}\longrightarrow
 \frac vD\frac{V_Z}{c+\delta}>\lambda\frac vD.
\]
One sufficiently large finite $K$ therefore restores every maintained
condition and increases the ratio by a factor greater than $\lambda$.

For any prescribed $C>0$, choose a positive integer $k$ with
$\lambda^k>4C^2$.  Perform $k$ such steps, choosing a finite batch
size separately at each step.  The final observation remains finite
and satisfies
\[
 v(F)/D>4C^2,\qquad \E|T_F|>\tfrac12\sqrt{v(F)}.
\]
Let $N$ now denote the total number of bits underlying the final
observation.  Its sign readout $f=\operatorname{sign}(T_F)$, with sign $1$ at zero,
has degree at most $D$.  Conditioning the bit sum on $F$ gives the
signed identity
\[
 \sum_i\widehat f(\{i\})=\E[fH_N]=\E|T_F|
 >\tfrac12\sqrt{v(F)}>C\sqrt D\ge C\sqrt{\deg(f)}.
\]
The positive signed correlation excludes constant $f$.  Arbitrarily
large target ratios follow from the same finite iteration.
\end{proof}


\subsection{Applying the criterion to the reporting event}

Apply the criterion first with $q=m+1$ and the reporting event $A$ of
Lemma~\ref{lem:gaussian}. That lemma gives positive probability,
conditional variance of the coordinate sum below one, and Gaussian-null
boundary. Identity~\eqref{eq:indicator} gives the required pointwise
expansion into functions of at most $m$ coordinates.
Lemma~\ref{lem:refinement} therefore supplies a finite partition with costs
$m$ and $m+1$ and with $V>c=m+1-\Pp(A)$.
Theorem~\ref{thm:mean-cost} gives the signed Boolean conclusion.
Appendix~\ref{app:weighted-mixture} supplies a second event with these
same four properties and uses the same refinement and amplification.
